\Page{Organiser son travail}{Parcours de travail}{parcours}
\thispagestyle{mpsiFront}
La progression du livre peut accompagner celle du professeur. Les outils de mesure, les incertitudes, le calcul et la programmation se travaillent dès les premières séances, au fil des thèmes.
\Sub{Choisir son point de départ}
\begingroup\renewcommand{\arraystretch}{1.28}
\begin{tabularx}{\linewidth}{@{}p{35mm}Y@{}}
\toprule
Parcours & Travail proposé dans le premier chapitre\\\midrule
Reprendre les bases & Diagnostic, cours, puis E01 à E09. Refaire les schémas et annoncer les hypothèses.\\[2mm]
Installer les méthodes & E10 à E18 : conjugaison, signes, foyers, projection, œil et appareil photographique.\\[2mm]
Gagner en autonomie & E19 à E24 : réglages, fibre, association de lentilles et incertitudes.\\[2mm]
Approfondir & Défis D1 à D3. Chercher d'abord seul ; consulter les indices progressivement.\\[2mm]
Préparer les devoirs & DM, puis DS de 2 h. Corriger avec le barème et reprendre les capacités fragiles.\\[2mm]
Se confronter à un support réel & Parcours R1 à R3 : retrouver le document original et traiter les questions ciblées.\\
\bottomrule
\end{tabularx}\endgroup
\Sub{Une séance de travail en quatre temps}
\begin{enumerate}
\item \textbf{Reformuler.} Dire avec ses propres mots ce qu'énonce une loi, avec ses conditions de validité et les unités des grandeurs.
\item \textbf{Chercher.} Dessiner la situation, nommer les inconnues et choisir une stratégie avant de remplacer les lettres par des nombres.
\item \textbf{Corriger.} Comparer les étapes du raisonnement ; noter la première erreur ou la première idée manquante.
\item \textbf{Réactiver.} Refaire la question quelques jours plus tard et expliquer la solution à voix haute, comme en interrogation orale.
\end{enumerate}
\begin{methode}{Un rythme à adapter}
Pour ce premier chapitre : consacrer une première étape au cours et aux constructions, une deuxième aux exercices et au TP, puis une troisième aux défis et évaluations. La compréhension et la capacité à refaire seul déterminent le passage à l'étape suivante.
\end{methode}
\source{Les difficultés et durées indiquées dans les exercices servent de repères. Le barème d'un DS est un outil pédagogique de ce livre, pas un barème officiel de concours.}
